\documentclass[12pt, titlepage]{article}
%\usepackage[norsk]{babel}
\usepackage[utf8]{inputenc}
\usepackage[T1]{fontenc}
\usepackage{minted, enumerate}
\usepackage[colorlinks = true,
linkcolor = blue,
urlcolor = blue]{hyperref}


\begin{document}

\title{Innlevering: Mappe}
\author{IIRA2001 - Haust 2024}
\date{}
\maketitle

Mappa består av dei to mappeoppgåvene vi har jobba med i haust.

\section*{Del I}

I fyrste del av mappa skulle de skrive 3 pythonprogram:
\begin{itemize}
   \item {Eit som brukar den logistiske likninga, til dømes ved populasjonsvekst}
   \item {Ein sparekalkulator, \textit{eller} eit program rundt databehandling av ei kundegruppe}
   \item {Ein simulering av kundeåtferd i butikk \textit{eller} simulering av marknadsdynamikk}
\end{itemize}

Sjå eigen oppgåvebeskriving på blackboard for meir detaljar om programma.

\section*{Del II}

I andre del av mappa skulle de lage 2 program som gjorde dataanalyse, brukte eit web-api eller begge delar.
Vi kravde at minst 1 av programma brukte data frå SSB eller Eurostat. Vi ser gjerne at de bruker eit web-api eller eit bibliotek til å hente dataa gjennom python i 1 (begge er også greitt) av programma.

Sjå eigen oppgåvebeskriving på blackboard for meir detaljar om programma.

\section*{Rapport}

\begin{itemize}
   \item{Til kvart pythonprogram skulle de også skrive ein liten rapport.}
   \item {Rapporten må innehalde ei beskriving av kva målet med programmet er, til dømes kva problemstilling du undersøkjer.}
   \item {Rapport og kodekommentar må til saman forklare korleis koden din fungerer.}
   \item {I tillegg treng du ein diskusjon eller analyse av resultatet frå programmet.}
\end{itemize}

Siste punkt er spesielt viktig og relevant for del 2 av mappa.

Kort fortalt treng vi ein liten rapport som beskriver kva målet med koden din var, og diskuterer kva resultatet av koden din vart, eller kva du fann ut.\\
Kor lang teksten er, avheng av programmet ditt.
Om du har laga ein sparekalkulator og brukar den til å finne ut kor mykje du må spare per månad for å kjøpe deg ein lamborghini i 30-årsgåve til deg sjølv, er det kanskje færre vurderingar og diskusjonar å skrive enn om du undersøkjer effekten straumprisen i Europa har på aksjekursen til Norsk Hydro.

~\\
De kan levere kode og rapport på fleire måtar, alt etter kva som passar deg best:
\begin{itemize}
   \item {Som 1 *.ipynb-fil (Jupyter-notebook-fil) per program med kodeceller for python og tekstceller for beskriving/resultat/diskusjon.}
   \item {Som 1 *.py-fil (pythonscript) med koden din og 1 PDF-fil med beskriving/resultat/diskusjon, per program.}
   \item {Som 1 *.py- eller *.ipynb-fil per program og ein samla rapport for heile mappa, eller 2 rapportar for del 1 og del 2 av mappa.}
\end{itemize}

Det er godkjent å ha fleire program i same *.py- eller *.ipynb-fil, men vi ser helst at dette blir unngått.
Leverer du rapporten i anna format enn PDF (t.d. Word), risikerer du at sensor ikkje klarer å lese den.

Hugs også å legge ved eventuelle data de har henta frå SSB eller andre kjelder (om de ikkje brukar web-api):

De har eigentleg fridom til å presentere mappa dykkar på den måten du meiner er best, men på eige ansvar:
\begin{itemize}
   \item {Vi må kunne køyre koden dykkar utan for mykje styr.}
   \item {Programma treng tekst som beskriv dei, og diskuterer eventuelle problemstillingar og resultat.}
\end{itemize}

\section*{KI-deklarasjonsskjema}

Til mappa skal det leggast ved signert og utfylt \href{https://i.ntnu.no/documents/portlet_file_entry/1305837853/NO_KI+-+deklareringsskjema+per+v%C3%A5r+2024.pdf/3b07aee3-d348-4360-e595-057ec480f1a3}{KI-deklarasjonsskjema}.

Det er rekna som juks å generere svar og kode ved hjelp av kunstig intelligens og levere den heilt eller delvis som eiga innlevering.
Ver nøye med å gjere greie for delar av kode og innlevering du ikkje har laga sjølv.

\section*{Inspera}

Det vert opna for innlevering på Inspera 9.12 klokka 09:00, og innleveringsfristen er 13.12 klokka 12:00 (ikkje midnatt).

Heile mappa med kode, rapport og KI-deklarasjonsskjema skal pakkast saman til ein zip-fil og lastast opp i Inspera.

Vi legg ut litt informasjon og lenker på Blackboard om korleis de pakkar saman filer til eit zip-arkiv, og kvar de finn Jupyter-notebook-filene dykkar.

\end{document}
